\documentclass[conference]{IEEEtran}
\usepackage{cite}
\usepackage{amsmath,amssymb,amsfonts}
\usepackage{graphicx}
\usepackage{array}
\usepackage{textcomp}
\usepackage{xcolor}
\usepackage{url}
\newcommand{\gf}{GenomeFlux}
\newcommand{\state}[1]{\texttt{#1}}
\begin{document}

\title{Why Are Some Genomes So Big and Others Very Small?\\
An Evidence-Gated, Phylogenetically Paired Framework for Genome-Size Evolution}

\author{}
\maketitle

\begin{abstract}
Nuclear genome size varies by orders of magnitude among eukaryotes without tracking organismal complexity or protein-coding gene number. In plants, large differences are often associated with repetitive DNA, transposable-element turnover, and polyploidy, but these gain processes alone cannot explain compact genomes, post-polyploid downsizing, or repeated bidirectional change in related lineages. This paper reframes the C-value question as a rate-balance inference problem: how do repeat-family amplification, whole-genome duplication (WGD), and DNA loss or compaction jointly explain present 1C genome size, and when can genome-size-associated phenotypes be interpreted as constraints rather than correlations? We propose GenomeFlux, a design-only comparative-genomics framework that fits gain-only, gain-plus-WGD, and bidirectional gain-loss models to phylogenetically structured, quality-gated data. The framework treats measurement provenance, repeat-annotation completeness, WGD dating, and sampling coverage as non-optional evidence conditions. It evaluates models by held-out-clade calibration rather than random species splitting, tests negative controls for annotation and phylogenetic confounding, and returns abstention when critical evidence is inadequate. The contribution is not a claim that a universal cause of genome size has been discovered. It is an auditable research design for deciding when a specific comparative dataset supports a conditional mechanism ranking, a lineage-specific explanation, or no directional conclusion.
\end{abstract}

\begin{IEEEkeywords}
genome size, C-value enigma, transposable elements, repeat DNA, polyploidy, genome contraction, comparative genomics, phylogenetic inference
\end{IEEEkeywords}

\section{Introduction}

Genome size is the amount of DNA in a haploid, unreplicated nuclear chromosome set (the 1C value). It is not a count of genes, a synonym for biological complexity, or a direct measure of adaptation. The empirical range is nevertheless extraordinary: Pellicer \emph{et al.} summarize more than 64,000-fold variation across eukaryotes and approximately 2,400-fold variation across land plants \cite{pellicer2018}. The extreme examples often used to motivate this question require careful wording. The microsporidian \emph{Encephalitozoon intestinalis} is among the smallest sequenced eukaryotic nuclear genomes, whereas the earlier flowering-plant record of \emph{Paris japonica} is no longer the largest known eukaryotic genome. A 2024 report identifies a 160.45 Gbp/1C genome in the fork fern \emph{Tmesipteris oblanceolata} \cite{fernandez2024}. These records reveal scale, not a single causal law.

The historical C-value enigma emerged because genome size does not scale simply with visible complexity. In many plant comparisons, noncoding and repetitive DNA account for a substantial portion of size differences \cite{pellicer2018,lee2014}. Transposable elements (TEs), particularly long terminal repeat retrotransposons, can multiply in lineage-specific bursts; whole-genome duplication (WGD) can instantly duplicate a genome; and small insertions, deletions, recombination, and repeat removal can subsequently reverse or erase part of that expansion. Genome size is therefore an evolutionary state variable produced by processes with different timescales.

The central scientific problem is not whether TEs, WGD, or deletion matter. Each has direct support in at least some systems. The harder problem is to infer their relative and time-dependent contributions without treating species as independent rows, without substituting assembly artifacts for biology, and without converting trait correlations into causal stories. For example, a large genome may coincide with cell or life-history differences, but that association can arise through common ancestry, correlated ecology, or reciprocal effects. Conversely, a compact genome may follow stringent DNA removal, weak repeat activity, or both.

This paper proposes \gf{}, a comparative framework that turns these qualitative alternatives into model comparisons. The framework is deliberately ambitious in scope--it seeks mechanism rankings that transfer beyond a single model species--but conservative in claim type. It produces conditional computational inferences only. It does not sample organisms, sequence DNA, alter genomes, observe evolution in real time, or state that any proposed mechanism has been demonstrated in a new lineage.

\section{Evidence Boundary and Research Gap}

\subsection{Repeat DNA is a mechanism, not a universal explanation}

TEs are endogenous mutators whose effects can range from parasitic propagation to neutral persistence or host co-option \cite{kidwell2001}. Their relevance to genome size is empirically concrete. In \emph{Gossypium}, differential, lineage-specific TE-family amplification explained substantial genome-size variation, with gypsy-like elements contributing major expansions in some lineages \cite{hawkins2006}. In the legume tribe Fabeae, repeat profiling of 23 species showed that differential repeat accumulation accounted for most observed genome-size differences, with one Ty3/gypsy lineage contributing a particularly large share \cite{macas2015}. The correct inference from these results is not that a total repeat percentage is sufficient. Distinct families can follow distinct histories, and an assemblage can be large because of a few successful lineages rather than broad TE diversity.

Large genome size is also not merely a TE count. Some genomes contain abundant satellites, segmental duplications, organellar insertions, or other noncoding components. Annotation methods, assembly continuity, and repeat libraries determine what is counted, particularly in repeat-rich genomes. Thus, a cross-species regression of 1C value on total annotated repeat fraction can create an apparently precise mechanism from inconsistent detection. A defensible analysis must preserve repeat-family resolution and encode annotation quality as part of the evidence, rather than as a post hoc nuisance covariate.

\subsection{Contraction and post-WGD change challenge one-directional models}

Genome expansion has an intuitive molecular narrative, but shrinkage is equally important. The carnivorous bladderwort \emph{Utricularia gibba} has an 82 Mb genome with a typical complement of plant genes despite evidence for repeated WGD events; much of its compactness reflects a strong reduction of non-genic DNA \cite{ibarra2013}. A comparative analysis in \emph{Genlisea} found divergent size evolution within one genus: an exceptionally small genome was consistent with retroelement silencing and deletion-biased repair, while a larger relative combined WGD and retrotransposition \cite{vu2015}. These examples establish that the same broad taxonomic group can move in opposite directions.

WGD is an initial condition rather than a permanent size verdict. Polyploid genomes undergo gene loss, rearrangement, epigenetic change, repeat turnover, and rediploidization \cite{soltis2003}. Consequently, a simple model that counts ancient duplication events will misattribute later DNA removal to the absence of duplication, or interpret a present large genome as an unmodified WGD product. The design therefore separates WGD timing and confidence from repeat gain and compaction proxies.

\subsection{Phenotypic relevance is plausible but causally underidentified}

Genome size has been associated with nuclear and cell properties, developmental timing, ecological distributions, and diversification. Gregory distinguishes mutation-pressure accounts from models in which bulk DNA has nucleotypic effects \cite{gregory2001}. Knight reviews evidence consistent with evolutionary and ecological constraints on large plant genomes \cite{knight2005}. Within a natural rotifer population, genome size covaries with body size, egg size, and embryonic development time \cite{stelzer2021}. These observations make trait constraints worth testing; they do not decide the direction of causation across eukaryotes.

\begin{table*}[t]
\caption{Evidence classes and permitted claims in the GenomeFlux proposal.}
\label{tab:evidence}
\centering
\footnotesize
\renewcommand{\arraystretch}{1.12}
\begin{tabular}{|p{0.17\textwidth}|p{0.20\textwidth}|p{0.24\textwidth}|p{0.22\textwidth}|}
\hline
\textbf{Evidence class} & \textbf{Representative sources} & \textbf{Permitted use} & \textbf{Use not permitted} \\
\hline
Genome-size diversity and records & \cite{pellicer2018,fernandez2024} & Define 1C value, diversity scale, and historical record boundary & Infer a universal biological maximum or mechanism \\
\hline
Repeat-family turnover & \cite{kidwell2001,hawkins2006,macas2015,lee2014} & Motivate family-resolved gain features and lineage comparisons & Equate total repeat fraction with an evolutionary rate \\
\hline
Compaction and bidirectional change & \cite{ibarra2013,vu2015} & Motivate separate loss/compaction features and paired contrasts & Assign a universal deletion mechanism to all small genomes \\
\hline
WGD and reorganization & \cite{soltis2003,ibarra2013} & Treat WGD as dated, uncertain historical input & Use WGD count as current genome size by itself \\
\hline
Trait associations & \cite{gregory2001,knight2005,stelzer2021} & Specify causal alternatives and confounder checks & State that a correlation proves adaptive genome size \\
\hline
\end{tabular}
\end{table*}

\subsection{Research gaps}

Five linked gaps follow from the evidence boundary. First, most comparisons are static: they do not jointly represent gain, loss, and dated WGD history, making a present 1C value compatible with multiple histories. Second, total repeat measures obscure whether a few TE families, broad repeat diversity, or unannotated components drive a contrast. Third, post-WGD dynamics confound current size with historical duplication. Fourth, trait constraints remain causally underidentified. Fifth, sparse taxonomic coverage and heterogeneous measurements can distort both the apparent range and inferred mechanisms. These gaps motivate a design that treats evidence quality and model abstention as scientific outputs rather than implementation failures.

\section{GenomeFlux Framework}

\subsection{Core representation}

For taxon or lineage $i$, \gf{} represents present log genome size as a conditional result of gain features $R_i$, WGD-history features $W_i$, loss or compaction features $D_i$, quality and measurement terms $Q_i$, and a phylogenetic effect $u_i$:

\begin{equation}
\log(C_i) = \alpha + \boldsymbol{\beta}_R^{\mathsf{T}}R_i + \boldsymbol{\beta}_W^{\mathsf{T}}W_i - \boldsymbol{\beta}_D^{\mathsf{T}}D_i + \boldsymbol{\beta}_Q^{\mathsf{T}}Q_i + u_i + \epsilon_i .
\label{eq:genomeflux}
\end{equation}

Here $C_i$ is a provenance-coded 1C estimate rather than an assembly span; $R_i$ may contain repeat-family abundance or age-distribution summaries; $W_i$ contains WGD timing and confidence; $D_i$ contains compaction or repeat-loss proxies; $Q_i$ records measurement, assembly, and annotation properties; and $u_i$ encodes phylogenetic covariance. Equation \eqref{eq:genomeflux} is a model specification, not an estimated mechanism or an observed data result. The sign convention makes explicit that putative loss features oppose DNA bulk, but the posterior sign and stability must be learned from qualified data.

The crucial design choice is comparison. Model A includes gain features and quality/phylogeny structure. Model B adds WGD history. Model C is the bidirectional \gf{} model in Eq. \eqref{eq:genomeflux}. A more complex model is not preferred because it tells a better story; it must improve held-out-clade predictive calibration and survive prespecified measurement, repeat-library, and WGD-date perturbations.

\begin{figure}[t]
\centering
\footnotesize
\begin{tabular}{c}
\textbf{GenomeFlux evidence-to-decision flow}\\[2pt]
\fbox{\parbox{0.86\columnwidth}{\centering Curated 1C values + phylogeny + measurement provenance}}\\[3pt]
$\downarrow$\\[-1pt]
\begin{tabular}{c@{\qquad}c@{\qquad}c}
\fbox{\parbox{0.22\columnwidth}{\centering Repeat-family gain}} &
\fbox{\parbox{0.22\columnwidth}{\centering WGD timing and rediploidization}} &
\fbox{\parbox{0.22\columnwidth}{\centering Loss and compaction proxies}}
\end{tabular}\\[3pt]
$\downarrow$\\[-1pt]
\fbox{\parbox{0.86\columnwidth}{\centering Phylogenetically paired models A/B/C + quality gates + negative controls}}\\[3pt]
$\downarrow$\\[-1pt]
\begin{tabular}{c@{\qquad}c}
\fbox{\parbox{0.38\columnwidth}{\centering Conditional mechanism ranking or lineage-specific result}} &
\fbox{\parbox{0.38\columnwidth}{\centering Critical evidence absent or unstable: evidence/model abstention}}
\end{tabular}
\end{tabular}
\caption{Conceptual \gf{} workflow. The editable diagram source is retained with the four-agent artifacts. The flow proposes a computational decision procedure; it reports no newly observed genome-size evolution.}
\label{fig:genomeflux}
\end{figure}

\subsection{Non-compensatory evidence gates}

Predictive scores cannot rescue a contrast whose input is not comparable. Let $E_{ik}$ be the evidence status for taxon $i$ in critical domain $k$, including 1C provenance, cytotype or ploidy coding, repeat annotation, WGD confidence, phylogenetic placement, and sampling independence. Let $T_{ik}$ be its transferability or comparability qualifier. Eligibility for a confirmatory contrast is

\begin{equation}
G_i = \prod_{k\in\mathcal{K}} \mathbb{I}\left(E_{ik}=\mathrm{sufficient}\right)\mathbb{I}\left(T_{ik}\geq\tau_k\right),
\label{eq:gate}
\end{equation}

where $\mathcal{K}$ is declared before fitting and $\tau_k$ is the threshold for domain $k$. If $G_i=0$, the record may be retained for descriptive coverage analysis but cannot create a confirmatory mechanism claim. This prevents rich assemblies from silently dominating an analysis because records from poorly sampled clades were treated as equivalent by default.

\section{Design-Only Methods}

\subsection{Data architecture and curation}

The proposed unit is a phylogenetically structured taxon record, not a field sample. Every record would store the 1C value and its publication or database source; measurement method and uncertainty; ploidy or cytotype; phylogenetic placement; assembly and repeat-library characteristics; repeat-family composition; WGD evidence and dating confidence; compaction or loss proxies; and permitted-claim ceiling. The first analysis would focus on land plants because the evidence base is substantial but heterogeneous. Selected outgroup clades may be held out as a transfer test only when their measurement and annotation fields meet the same gates.

No missing value is automatically converted to biological absence. An unknown TE family profile is not a low TE profile; an incomplete WGD history is not evidence of no duplication; and a genome assembly span is not necessarily a nuclear 1C value. The data dictionary therefore records whether a field is observed, derived, assumed for sensitivity analysis, or unavailable. Records failing a critical field receive an exclusion reason or an abstention state rather than an imputed confirmatory label.

\subsection{Paired validation and negative controls}

Random train-test splitting would place close relatives on both sides of the test and exaggerate generalization. \gf{} instead withholds related clades, reconstructs phylogenetically paired contrasts, and evaluates whether a model trained on other lineages predicts the held-out contrast. For a model $m$ under scenario $s$, the comparison summary combines held-out calibration $L_{ms}$, residual phylogenetic structure $P_{ms}$, and a declared complexity penalty $K_m$:

\begin{equation}
\mathcal{V}_{ms}=L_{ms}-\lambda_P P_{ms}-\lambda_K K_m,
\label{eq:validation}
\end{equation}

where $\lambda_P$ and $\lambda_K$ are fixed before analysis. Equation \eqref{eq:validation} is a transparent selection aid, not a claim that one scalar captures scientific truth. The report must show its components, uncertainty intervals, and the effect of leaving out each target clade.

Negative controls are mandatory. Shuffling repeat-family labels tests whether a model gains apparent skill from taxonomic grouping rather than family-specific biology. Permuting WGD dates tests whether historical timing contains information beyond a generic lineage label. Quality-matched subsets test whether assembly or annotation differences create the result. A trait analysis adds a separate causal-alternative set, including phylogeny and declared life-history covariates; it is never folded into the mechanistic genome-size model as an unexamined “benefit” score.

\begin{table}[t]
\caption{Predeclared models and primary falsification logic.}
\label{tab:models}
\centering
\scriptsize
\renewcommand{\arraystretch}{1.1}
\begin{tabular}{|p{0.16\columnwidth}|p{0.26\columnwidth}|p{0.19\columnwidth}|p{0.15\columnwidth}|}
\hline
\textbf{Model} & \textbf{Terms} & \textbf{Question addressed} & \textbf{Failure condition} \\
\hline
A: Gain-only & Repeat-family features, quality, phylogeny & Do gain features alone transfer across clades? & No calibration advantage over baseline \\
\hline
B: Gain + WGD & A plus dated WGD history & Does WGD add predictive information beyond repeat features? & Timing term unstable or no held-out gain \\
\hline
C: GenomeFlux & B plus compaction/loss features and hierarchical clade effects & Does bidirectional history improve robust explanation? & A/B equal or better after complexity penalty \\
\hline
Trait extension & Genome size, trait data, phylogeny, causal alternatives & Is a trait association robust enough to motivate follow-up? & Association vanishes or remains non-identifiable \\
\hline
\end{tabular}
\end{table}

\subsection{Scenario matrix and reproducibility}

Six scenarios make sensitivity explicit. The baseline uses a preregistered curated record set. A conservative-measurement scenario removes weakly documented 1C values. A repeat-library scenario changes the annotation library or classification rule. A WGD-date scenario propagates alternative dating and confidence states. A leave-one-clade-out scenario measures transfer rather than within-clade memorization. A trait-confounding scenario tests whether any phenotypic association survives declared controls. Results are summarized by the minimum supported conclusion across eligible scenarios, not by the most favorable single run.

Reproducibility requires a versioned evidence registry. Every published comparison should retain input card identifiers, filtering decisions, model code version, phylogeny version, scenario register, gate status, uncertainty class, and human-review notes. A later improved annotation creates a new analysis version; it does not silently rewrite a prior claim. This distinction matters because progress in sequencing can change the observability of repeat-rich genomes without changing their evolutionary history.

\section{Conditional Results Logic and Limitations}

The study will not report observed TE bursts, deletion rates, WGD events, trait effects, or evolved genome-size changes. It will report conditional model behavior on a future qualified dataset. Four branches are informative. First, Model C may outperform A and B across independent clades; the permitted statement would be that a combined gain-loss representation improves conditional comparative calibration. Second, different clades may favor different histories; the permitted statement would be that mechanisms are lineage-specific. Third, a simpler model may match C; the added mechanism would not be justified for that dataset. Fourth, a quality gate may fail; the valid result is \state{EVIDENCE\_OR\_MODEL\_ABSTAIN} plus a precise evidence request.

Several limitations remain material. The literature synthesis is targeted, not systematic. Published C-values are not uniformly comparable. Repeat-rich genomes may be incompletely assembled or annotated, and the availability of high-quality data is taxonomically uneven. Loss proxies are not direct rates unless their mapping has independently been validated. WGD dates contain uncertainty and can be correlated with other lineage attributes. Finally, a phylogenetically robust trait association still does not establish whether genome size causes the trait, the trait changes selection on DNA amount, or both respond to another variable.

Human review is required before a biological conclusion is issued. Genome-size database expertise is needed for measurement provenance; comparative-genomics expertise is needed for repeat annotation; evolutionary-methods expertise is needed for phylogenetic structure; and evolutionary ecology expertise is needed for trait interpretation. The proposal neither replaces these judgments nor turns weakly represented taxa into evidence-free placeholders.

\section{Discussion and Conclusion}

The best current answer to why some genomes are enormous and others compact is plural rather than evasive. Genome size reflects a balance of DNA accumulation, DNA removal, genome-wide events, and lineage-specific context. Repetitive DNA and TEs are often central, but their family composition and activity vary. WGD can expand a genome instantly, but rediploidization and compaction can make ancient duplication an unreliable predictor of present bulk. Genome size can have biological consequences, yet the existence, direction, and strength of those consequences must be evaluated separately from the molecular history that produced the DNA amount.

\gf{} makes this answer operational. It rejects a gain-only default, prevents WGD from standing in for all history, and prohibits predictive fit from compensating for non-comparable measurement or annotation evidence. Its contribution is a falsifiable standard: a proposed mechanism must add out-of-clade explanatory value, survive prespecified perturbations, and remain interpretable after phylogenetic and quality controls. A failed model is useful because it identifies whether the missing piece is repeat-family resolution, loss-process evidence, WGD dating, taxon sampling, or trait data.

This is a proposal for a research program, not a claim that the C-value enigma has been solved. The framework will be most successful if it produces a map of heterogeneous mechanisms and explicit abstentions rather than a deceptively universal ranking. By making data provenance, model alternatives, negative controls, and limits of transferability visible, it creates a rigorous path from the observation of extreme genomes to testable evolutionary explanations.

\appendices
\section{Scenario Register and Claim-Validation Contract}

Before fitting, a future \gf{} project must publish the eligible taxon set, phylogeny source, 1C inclusion rule, repeat-class taxonomy, WGD coding rule, loss/compaction proxy definitions, quality thresholds, transformations, clade hold-outs, and all comparison metrics. The registry separates observed input values from literature-derived annotations and from values used only in sensitivity scenarios. No record may move from “unknown” to “absent” merely because a model requires a complete matrix.

The registry must also define the expected failure signatures. A mechanism result is unstable if its sign changes across reasonable repeat-library or measurement filters, if it is supported by one unreplicated clade only, if shuffled labels retain the same explanatory gain, or if the result vanishes under quality-matched subsets. A WGD result is unstable if its direction depends on a single dating scheme. A trait association is non-identifiable if different causal alternatives fit the available data equally well. These conditions are scientific findings about the limits of an inference, not optional caveats.

Each final sentence that interprets a model receives a claim ledger entry. The entry links to the relevant source cards, input-version identifier, model, gate status, held-out-clade result, uncertainty qualifier, and reviewer note. Allowed language ranges from “descriptive diversity statement” through “conditional comparative support.” Prohibited language includes “proved universal cause,” “observed evolutionary outcome,” and any implication that a public-data model has manipulated or measured a living lineage. The ledger preserves negative and abstention outcomes so that later data can refine, rather than erase, the record of uncertainty.

\begin{table*}[t]
\caption{Scenario register for a future GenomeFlux implementation. All entries are proposed analyses, not completed runs.}
\label{tab:scenarios}
\centering
\footnotesize
\renewcommand{\arraystretch}{1.14}
\begin{tabular}{|p{0.16\textwidth}|p{0.24\textwidth}|p{0.24\textwidth}|p{0.22\textwidth}|}
\hline
\textbf{Scenario} & \textbf{What changes} & \textbf{Failure to detect} & \textbf{Permitted response} \\
\hline
Conservative measurement & Exclude weakly documented 1C values and ambiguous cytotypes & Apparent size contrast depends on provenance-poor records & Downgrade to descriptive coverage or abstain \\
\hline
Repeat-library alternative & Change repeat classification or annotation library & Claimed family effect is a library artifact & Report annotation sensitivity; do not rank mechanism \\
\hline
WGD-date uncertainty & Propagate age and confidence alternatives & WGD term depends on one historical reconstruction & Retain WGD as uncertain input only \\
\hline
Leave-one-clade-out & Withhold related lineages as a test unit & Random row split has overstated transfer & Report lineage-specific result or no transfer \\
\hline
Quality-matched subset & Match assembly and annotation characteristics & Data-quality structure carries apparent mechanism & Apply evidence gate or abstain \\
\hline
Trait-confounding set & Add phylogeny and declared covariates & Trait link is common ancestry or correlated ecology & Report association only; prohibit causal claim \\
\hline
\end{tabular}
\end{table*}

\section{Human Review and Evidence-Coverage Record}

The analysis remains multidisciplinary even though it does not perform wet-lab work. A curator verifies C-value units, standards, cytotypes, and database provenance. A repeat-biology reviewer assesses whether family labels and activity proxies are comparable. A phylogeneticist reviews branch lengths, pairing logic, and model covariance. A genome-evolution specialist evaluates WGD coding and whether compaction proxies justify their stated interpretation. An evolutionary ecologist reviews any trait-extension claim. These roles should record disagreement rather than average it into a numerical weight.

Coverage is also an equity and reliability issue. Species with polished assemblies and comprehensive annotations will be overrepresented unless the project reports the taxa and genome architectures it cannot currently assess. A high abstention rate in underrepresented clades is not a failure of those lineages; it is evidence that current infrastructure cannot support the same mechanism claim. The appropriate next step is targeted data improvement and expert curation, not a global extrapolation from model-rich species.

\section{Model-Interpretation and Escalation Ledger}

The framework is intentionally designed to make a strong comparison without collapsing it into a universal claim. A future report must distinguish three levels of successful output. At the first level, a model can establish a descriptive pattern, such as a repeat-family contrast within a qualified clade. At the second, it can establish conditional comparative support when the pattern survives held-out-clade validation and the stated data-quality scenarios. At the third, it can motivate a more direct mechanistic study only after the gain, loss, and historical alternatives have been separately examined. No level permits a conclusion to skip directly from “predictive coefficient” to “the cause of genome size.”

The escalation ledger records what would be needed to strengthen each type of claim. An unstable repeat-family term calls for repeat-library reconciliation, targeted high-contiguity assemblies, or independently curated repeat annotations. An unstable compaction term calls for better time-resolved loss evidence rather than a larger regression alone. A WGD result that depends on one reconstruction calls for dated phylogenomic review. A trait result that survives a simple correlation but fails the causal-alternative set calls for a design that can distinguish developmental cost, ecological selection, and common ancestry. In this way, a negative model result becomes an actionable map of experimental and curation priorities.

The ledger must preserve counterexamples. Compact \emph{Utricularia} after repeated WGD and divergent \emph{Genlisea} histories are not inconvenient exceptions to remove; they are the observations that falsify a monotonic expansion story. Likewise, a large genome with no current signal of an active TE burst should not be forced into a recent-gain category merely because it is large. The historical state may be old, the relevant repeat family may be missing from the library, or removal may be weak rather than gain unusually strong. The planned model comparison is valuable only if its outputs retain those alternatives.

Finally, the report should make the difference between discovery and deployment explicit. GenomeFlux can guide a comparative research program by prioritizing which clades, annotations, and hypotheses are most informative. It cannot decide conservation status, crop value, organismal fitness, or any intervention on the basis of C-value alone. The framework's ambition is epistemic: it replaces a broad riddle with a sequence of reproducible discriminations among mechanisms, each of which can be strengthened, overturned, or appropriately withheld as evidence improves.

\section{Minimum Data Record and Benchmark Protocol}

The proposed study requires a machine-readable record for every taxon before any model is compared. The point is not bureaucratic completeness. Genome-size research combines measurements, assemblies, repeat annotations, and evolutionary reconstructions that were generated for different purposes and at different resolutions. If the record does not state which kind of evidence supports a field, an apparent biological contrast can be an artifact of substitution: assembly span for C-value, uncurated annotation for family abundance, or inferred absence for missing data. The record turns those substitutions into detectable validation failures.

\begin{table*}[t]
\caption{Minimum evidence record for each GenomeFlux taxon. The fields define a future reproducible data contract, not a completed database.}
\label{tab:record}
\centering
\footnotesize
\renewcommand{\arraystretch}{1.14}
\begin{tabular}{|p{0.16\textwidth}|p{0.20\textwidth}|p{0.24\textwidth}|p{0.23\textwidth}|}
\hline
\textbf{Record block} & \textbf{Required fields} & \textbf{Quality question} & \textbf{Failure response} \\
\hline
Nuclear DNA measurement & 1C value, unit, method, standard, cytotype, source, uncertainty & Is this a comparable nuclear quantity rather than an assembly span or ambiguous estimate? & Exclude from confirmatory contrast; retain coverage note \\
\hline
Phylogenetic context & Accepted name, lineage, tree version, branch information, pairing eligibility & Are the target and comparator placed consistently and independently enough for the split? & Do not use random row split as a substitute \\
\hline
Repeat architecture & Library version, family class, abundance basis, assembly or read context, annotation quality & Can family abundance be compared without silent detection bias? & Run alternative library scenario or abstain \\
\hline
Genome-wide history & WGD evidence, timing interval, confidence, rediploidization notes & Is WGD historical input distinguished from present bulk DNA? & Treat as uncertain covariate or remove WGD claim \\
\hline
Loss and compaction & Intergenic context, repeat absence pattern, source, proxy mapping, uncertainty & Does the proxy support a loss/compaction interpretation at the stated ceiling? & Downgrade to descriptive pattern \\
\hline
Claim audit & Gate status, scenario, model version, hold-out result, reviewer note & Can a reader reconstruct why a conclusion was permitted? & Emit evidence/model abstention \\
\hline
\end{tabular}
\end{table*}

Benchmarking begins with a baseline that is intentionally modest: can a model predict a held-out clade's genome-size contrast better than a phylogeny-and-quality baseline? The gain-only, gain-plus-WGD, and bidirectional models are then compared on identical clade splits and the same eligible records. A model that gains accuracy only by seeing near relatives has not learned a transferable mechanism. A model that is slightly less accurate but consistently identifies an inadequate record may be more scientifically useful than one that hides incompatibility through imputation. Thus, the benchmark reports calibration, residual structure, complexity, and abstention rate together.

The protocol should also publish a challenge set containing deliberately difficult cases. One class includes compact genomes after WGD, such as \emph{Utricularia}, which tests whether duplication is being mistaken for current bulk. A second includes related lineages with opposite repeat trajectories, such as \emph{Genlisea}, which tests whether the model can represent bidirectionality. A third includes repeat-rich or poorly assembled genomes, which tests whether data-quality gates prevent spurious family effects. A fourth includes trait-rich lineages, which tests whether a phenotypic correlation is mistakenly absorbed into the molecular mechanism. These are not demonstrations of expected performance; they are preregistered attempts to break an overconfident model.

The final benchmark artifact should include both a human-readable report and a compact result table. For each clade and scenario it must state the model ranking, uncertainty, critical evidence status, reason for any abstention, and the strongest permissible sentence. A later study may revise the table when new measurements or annotations are added, but it must preserve the older version and document what changed. This makes improvement cumulative: new data can resolve a named uncertainty without retroactively converting a proposal-level mechanism into a historical observation.

The benchmark should additionally report where the model fails gracefully. A practical implementation needs an error taxonomy that separates data absence, metadata ambiguity, model instability, and biological disagreement. Data absence means that a required field was never measured or made available; the appropriate output is a targeted acquisition request. Metadata ambiguity means that values exist but their unit, cytotype, source, or annotation basis cannot be reconciled; the appropriate output is exclusion from the confirmatory cohort. Model instability means that a coefficient, prediction, or mechanism ranking changes under a preregistered reasonable scenario; the appropriate output is an uncertainty statement, not selection of the preferred run. Biological disagreement means that independently qualified clades support different histories; the appropriate output is a lineage-specific result rather than forced consensus.

This taxonomy is especially valuable for very large and very small genomes because extremes attract overinterpretation. A giant genome can be technically difficult to assemble and annotate, while a compact genome can make missing sequence look like genuine compaction. A quality-aware benchmark asks whether the same qualitative conclusion survives when evidence classes are matched, not whether the most data-rich species yields the most striking coefficient. It also makes the next investment visible: a clade with uncertain C-values needs measurement curation; a clade with incompatible repeat annotations needs a harmonized library; a clade with rival WGD histories needs phylogenomic resolution; and a clade with stable molecular associations but unclear phenotype needs a distinct ecological design.

An implementation report should publish a compact failure matrix beside its model scores. For each rejected or abstained comparison, the matrix records the first failed gate, whether the failure is remediable, which expert role should review it, and which new information could change the decision. This prevents a future analyst from treating abstention as a negative biological result or from silently rerunning the model until a preferred answer appears. Reproducibility is therefore not merely code release. It is the ability to reconstruct why a promising-looking comparison was not allowed to become a mechanistic conclusion.

\IEEEtriggeratref{4}
\begin{thebibliography}{00}
\bibitem{pellicer2018} J. Pellicer, O. Hidalgo, S. Dodsworth, and I. J. Leitch, ``Genome Size Diversity and Its Impact on the Evolution of Land Plants,'' \emph{Genes}, vol. 9, no. 2, p. 88, 2018, doi: 10.3390/genes9020088.
\bibitem{fernandez2024} P. Fernandez \emph{et al.}, ``A 160 Gbp fork fern genome shatters size record for eukaryotes,'' \emph{iScience}, vol. 27, p. 109889, 2024, doi: 10.1016/j.isci.2024.109889.
\bibitem{kidwell2001} M. G. Kidwell and D. Lisch, ``Perspective: Transposable Elements, Parasitic DNA, and Genome Evolution,'' \emph{Evolution}, vol. 55, no. 1, pp. 1--24, 2001, doi: 10.1111/j.0014-3820.2001.tb01268.x.
\bibitem{hawkins2006} J. S. Hawkins \emph{et al.}, ``Differential lineage-specific amplification of transposable elements is responsible for genome size variation in \emph{Gossypium},'' \emph{Genome Research}, vol. 16, no. 10, pp. 1252--1261, 2006, doi: 10.1101/gr.5282906.
\bibitem{macas2015} J. Macas \emph{et al.}, ``In Depth Characterization of Repetitive DNA in 23 Plant Genomes Reveals Sources of Genome Size Variation in the Legume Tribe Fabeae,'' \emph{PLoS ONE}, vol. 10, p. e0143424, 2015, doi: 10.1371/journal.pone.0143424.
\bibitem{ibarra2013} E. Ibarra-Laclette \emph{et al.}, ``Architecture and evolution of a minute plant genome,'' \emph{Nature}, vol. 498, pp. 94--98, 2013, doi: 10.1038/nature12132.
\bibitem{vu2015} G. T. H. Vu \emph{et al.}, ``Comparative Genome Analysis Reveals Divergent Genome Size Evolution in a Carnivorous Plant Genus,'' \emph{The Plant Genome}, vol. 8, p. eplantgenome2015.04.0021, 2015, doi: 10.3835/plantgenome2015.04.0021.
\bibitem{soltis2003} P. S. Soltis, P. S. Soltis, and J. A. Tate, ``Advances in the study of polyploidy since \emph{Plant speciation},'' \emph{New Phytologist}, vol. 161, no. 1, pp. 173--191, 2003, doi: 10.1046/j.1469-8137.2003.00948.x.
\bibitem{gregory2001} T. R. Gregory, ``Coincidence, coevolution, or causation? DNA content, cell size, and the C-value enigma,'' \emph{Biological Reviews}, vol. 76, no. 1, pp. 65--101, 2001, doi: 10.1111/j.1469-185x.2000.tb00059.x.
\bibitem{knight2005} C. A. Knight, ``The Large Genome Constraint Hypothesis: Evolution, Ecology and Phenotype,'' \emph{Annals of Botany}, vol. 95, no. 1, pp. 177--190, 2005, doi: 10.1093/aob/mci011.
\bibitem{stelzer2021} C.-P. Stelzer, M. Pichler, and A. Hatheuer, ``Linking genome size variation to population phenotypic variation within the rotifer, \emph{Brachionus asplanchnoidis},'' \emph{Communications Biology}, vol. 4, p. 596, 2021, doi: 10.1038/s42003-021-02131-z.
\bibitem{lee2014} S.-I. Lee and N.-S. Kim, ``Transposable Elements and Genome Size Variations in Plants,'' \emph{Genomics \& Informatics}, vol. 12, no. 3, p. 87, 2014, doi: 10.5808/gi.2014.12.3.87.
\end{thebibliography}

\end{document}
